\documentclass[11pt]{article}
% Shared preamble for the PNPInverse onboarding docs.
% Each document does % Shared preamble for the PNPInverse onboarding docs.
% Each document does % Shared preamble for the PNPInverse onboarding docs.
% Each document does \input{preamble} after \documentclass.

\usepackage[utf8]{inputenc}
\usepackage[T1]{fontenc}
\usepackage{lmodern}
\usepackage[margin=1in]{geometry}
\usepackage{amsmath, amssymb, amsthm}
\usepackage{mathtools}
\usepackage{empheq}
\usepackage{siunitx}
\usepackage{booktabs}
\usepackage{array}
\usepackage{enumitem}
\usepackage{xcolor}
\usepackage{framed}
\usepackage{microtype}
\usepackage{listings}
\usepackage{tikz}
\usetikzlibrary{arrows.meta, positioning, decorations.pathreplacing, calligraphy}
\usepackage[hidelinks]{hyperref}

% --- branding -------------------------------------------------------------
\definecolor{pnpblue}{RGB}{30,70,120}
\definecolor{pnpgray}{RGB}{90,90,90}
\definecolor{calloutbg}{RGB}{245,245,235}
\definecolor{shadecolor}{RGB}{245,245,235}

% --- callouts (uses framed.sty's shaded environment) ----------------------
\newenvironment{callout}[1]
  {\begin{shaded}\noindent\textbf{\color{pnpblue}#1}\par\smallskip\ignorespaces}
  {\end{shaded}}

\newcommand{\concept}[1]{\textbf{\color{pnpblue}#1}}
\newcommand{\warn}[1]{\textbf{\color{red!70!black}#1}}

% --- code listings --------------------------------------------------------
\lstset{
  basicstyle=\ttfamily\small,
  breaklines=true,
  columns=fullflexible,
  keepspaces=true,
  showstringspaces=false,
  frame=single,
  framesep=4pt,
  backgroundcolor=\color{calloutbg},
  rulecolor=\color{pnpgray!40},
}

% --- math shortcuts -------------------------------------------------------
\newcommand{\R}{\mathbb{R}}
\newcommand{\diff}{\mathrm{d}}
\newcommand{\eqd}{\overset{!}{=}}
\newcommand{\grad}{\nabla}
\newcommand{\divg}{\nabla\!\cdot}
\newcommand{\eqdef}{\;\coloneqq\;}
\newcommand{\evalat}[2]{\left.#1\right\rvert_{#2}}
\newcommand{\set}[1]{\{#1\}}

% common chemistry / physics symbols
\newcommand{\Hyd}{\mathrm{H}}
\newcommand{\Oxy}{\mathrm{O}}
\newcommand{\HtwoO}{\mathrm{H_2O}}
\newcommand{\HtwoOtwo}{\mathrm{H_2O_2}}
\newcommand{\Otwo}{\mathrm{O_2}}
\newcommand{\Hp}{\mathrm{H^+}}
\newcommand{\OHm}{\mathrm{OH^-}}
\newcommand{\Kp}{\mathrm{K^+}}
\newcommand{\Csp}{\mathrm{Cs^+}}
\newcommand{\SOfour}{\mathrm{SO_4^{2-}}}
\newcommand{\Vrhe}{V_{\mathrm{RHE}}}
\newcommand{\eo}{E^{\circ}}
\newcommand{\Eeq}{E_{\mathrm{eq}}}
\newcommand{\kzero}{k_0}

\newcommand{\onboardingfooter}{%
  \vfill\noindent\rule{\linewidth}{0.4pt}\par
  {\small\color{pnpgray} PNPInverse intern onboarding,
  compiled \today. Source under \texttt{docs/onboarding/}.}\par}
 after \documentclass.

\usepackage[utf8]{inputenc}
\usepackage[T1]{fontenc}
\usepackage{lmodern}
\usepackage[margin=1in]{geometry}
\usepackage{amsmath, amssymb, amsthm}
\usepackage{mathtools}
\usepackage{empheq}
\usepackage{siunitx}
\usepackage{booktabs}
\usepackage{array}
\usepackage{enumitem}
\usepackage{xcolor}
\usepackage{framed}
\usepackage{microtype}
\usepackage{listings}
\usepackage{tikz}
\usetikzlibrary{arrows.meta, positioning, decorations.pathreplacing, calligraphy}
\usepackage[hidelinks]{hyperref}

% --- branding -------------------------------------------------------------
\definecolor{pnpblue}{RGB}{30,70,120}
\definecolor{pnpgray}{RGB}{90,90,90}
\definecolor{calloutbg}{RGB}{245,245,235}
\definecolor{shadecolor}{RGB}{245,245,235}

% --- callouts (uses framed.sty's shaded environment) ----------------------
\newenvironment{callout}[1]
  {\begin{shaded}\noindent\textbf{\color{pnpblue}#1}\par\smallskip\ignorespaces}
  {\end{shaded}}

\newcommand{\concept}[1]{\textbf{\color{pnpblue}#1}}
\newcommand{\warn}[1]{\textbf{\color{red!70!black}#1}}

% --- code listings --------------------------------------------------------
\lstset{
  basicstyle=\ttfamily\small,
  breaklines=true,
  columns=fullflexible,
  keepspaces=true,
  showstringspaces=false,
  frame=single,
  framesep=4pt,
  backgroundcolor=\color{calloutbg},
  rulecolor=\color{pnpgray!40},
}

% --- math shortcuts -------------------------------------------------------
\newcommand{\R}{\mathbb{R}}
\newcommand{\diff}{\mathrm{d}}
\newcommand{\eqd}{\overset{!}{=}}
\newcommand{\grad}{\nabla}
\newcommand{\divg}{\nabla\!\cdot}
\newcommand{\eqdef}{\;\coloneqq\;}
\newcommand{\evalat}[2]{\left.#1\right\rvert_{#2}}
\newcommand{\set}[1]{\{#1\}}

% common chemistry / physics symbols
\newcommand{\Hyd}{\mathrm{H}}
\newcommand{\Oxy}{\mathrm{O}}
\newcommand{\HtwoO}{\mathrm{H_2O}}
\newcommand{\HtwoOtwo}{\mathrm{H_2O_2}}
\newcommand{\Otwo}{\mathrm{O_2}}
\newcommand{\Hp}{\mathrm{H^+}}
\newcommand{\OHm}{\mathrm{OH^-}}
\newcommand{\Kp}{\mathrm{K^+}}
\newcommand{\Csp}{\mathrm{Cs^+}}
\newcommand{\SOfour}{\mathrm{SO_4^{2-}}}
\newcommand{\Vrhe}{V_{\mathrm{RHE}}}
\newcommand{\eo}{E^{\circ}}
\newcommand{\Eeq}{E_{\mathrm{eq}}}
\newcommand{\kzero}{k_0}

\newcommand{\onboardingfooter}{%
  \vfill\noindent\rule{\linewidth}{0.4pt}\par
  {\small\color{pnpgray} PNPInverse intern onboarding,
  compiled \today. Source under \texttt{docs/onboarding/}.}\par}
 after \documentclass.

\usepackage[utf8]{inputenc}
\usepackage[T1]{fontenc}
\usepackage{lmodern}
\usepackage[margin=1in]{geometry}
\usepackage{amsmath, amssymb, amsthm}
\usepackage{mathtools}
\usepackage{empheq}
\usepackage{siunitx}
\usepackage{booktabs}
\usepackage{array}
\usepackage{enumitem}
\usepackage{xcolor}
\usepackage{framed}
\usepackage{microtype}
\usepackage{listings}
\usepackage{tikz}
\usetikzlibrary{arrows.meta, positioning, decorations.pathreplacing, calligraphy}
\usepackage[hidelinks]{hyperref}

% --- branding -------------------------------------------------------------
\definecolor{pnpblue}{RGB}{30,70,120}
\definecolor{pnpgray}{RGB}{90,90,90}
\definecolor{calloutbg}{RGB}{245,245,235}
\definecolor{shadecolor}{RGB}{245,245,235}

% --- callouts (uses framed.sty's shaded environment) ----------------------
\newenvironment{callout}[1]
  {\begin{shaded}\noindent\textbf{\color{pnpblue}#1}\par\smallskip\ignorespaces}
  {\end{shaded}}

\newcommand{\concept}[1]{\textbf{\color{pnpblue}#1}}
\newcommand{\warn}[1]{\textbf{\color{red!70!black}#1}}

% --- code listings --------------------------------------------------------
\lstset{
  basicstyle=\ttfamily\small,
  breaklines=true,
  columns=fullflexible,
  keepspaces=true,
  showstringspaces=false,
  frame=single,
  framesep=4pt,
  backgroundcolor=\color{calloutbg},
  rulecolor=\color{pnpgray!40},
}

% --- math shortcuts -------------------------------------------------------
\newcommand{\R}{\mathbb{R}}
\newcommand{\diff}{\mathrm{d}}
\newcommand{\eqd}{\overset{!}{=}}
\newcommand{\grad}{\nabla}
\newcommand{\divg}{\nabla\!\cdot}
\newcommand{\eqdef}{\;\coloneqq\;}
\newcommand{\evalat}[2]{\left.#1\right\rvert_{#2}}
\newcommand{\set}[1]{\{#1\}}

% common chemistry / physics symbols
\newcommand{\Hyd}{\mathrm{H}}
\newcommand{\Oxy}{\mathrm{O}}
\newcommand{\HtwoO}{\mathrm{H_2O}}
\newcommand{\HtwoOtwo}{\mathrm{H_2O_2}}
\newcommand{\Otwo}{\mathrm{O_2}}
\newcommand{\Hp}{\mathrm{H^+}}
\newcommand{\OHm}{\mathrm{OH^-}}
\newcommand{\Kp}{\mathrm{K^+}}
\newcommand{\Csp}{\mathrm{Cs^+}}
\newcommand{\SOfour}{\mathrm{SO_4^{2-}}}
\newcommand{\Vrhe}{V_{\mathrm{RHE}}}
\newcommand{\eo}{E^{\circ}}
\newcommand{\Eeq}{E_{\mathrm{eq}}}
\newcommand{\kzero}{k_0}

\newcommand{\onboardingfooter}{%
  \vfill\noindent\rule{\linewidth}{0.4pt}\par
  {\small\color{pnpgray} PNPInverse intern onboarding,
  compiled \today. Source under \texttt{docs/onboarding/}.}\par}


\title{\bfseries Numerics \& Codebase Tour \\
       \large from Newton's method to the production solver}
\author{Onboarding doc 03 of 06}
\date{\today}

\begin{document}
\maketitle

\begin{callout}{What this document is for}
A guided tour of how the equations from doc 02 turn into running
code. Covers the finite-element library (Firedrake), the nonlinear
solver strategies (Newton + continuation + Picard), and the layout
of the actual repository. By the end you should be able to find
where any reaction term, IC seed, or solver flag lives.
\end{callout}

\tableofcontents
\bigskip\hrule\bigskip

% =========================================================================
\section{The pipeline at a glance}

\begin{center}
\begin{tikzpicture}[
  node distance=8mm and 14mm,
  box/.style={draw, rounded corners=3pt, minimum height=8mm,
              minimum width=32mm, align=center, font=\small},
  arr/.style={-{Latex[length=2mm]}}
]
\node[box, fill=blue!10]   (driver) {driver script\\(\texttt{scripts/studies/...py})};
\node[box, right=of driver, fill=blue!10] (factory) {\texttt{make\_bv\_solver\_params}\\\texttt{scripts/\_bv\_common.py}};
\node[box, right=of factory, fill=calloutbg] (orch) {orchestrator\\\texttt{anchor\_continuation.py}};
\node[box, below=of orch, fill=orange!20] (forms) {residual builder\\\texttt{forms\_logc\_muh.py}};
\node[box, left=of forms, fill=orange!20] (newton) {Firedrake\\Newton + PETSc};
\node[box, left=of newton, fill=green!15] (results) {\texttt{StudyResults/...}\\(JSON, plots)};

\draw[arr] (driver)  -- (factory);
\draw[arr] (factory) -- (orch);
\draw[arr] (orch)    -- (forms);
\draw[arr] (forms)   -- (newton);
\draw[arr] (newton)  -- (results);
\draw[arr] (newton.north) .. controls +(0,1) and +(0,-1) .. (orch.south west) node[midway, fill=white, font=\scriptsize] {next $V$};
\end{tikzpicture}
\end{center}

A driver script picks bulk concentrations, kinetic constants, and a
voltage grid; it asks the factory for a packed
\texttt{solver\_params} object; it hands those to the orchestrator,
which loops over voltages and continuation rungs; for each rung the
orchestrator builds the FEM residual via the forms file and asks
Firedrake to Newton-solve it. Converged solutions get post-processed
into a JSON of observables and saved under \texttt{StudyResults/}.

% =========================================================================
\section{Finite elements in 60 seconds}

You saw the weak form briefly in doc 02 \S 9. The
\concept{finite-element method (FEM)} approximates each unknown
function $u(x)$ as
\[
  u_h(x) = \sum_{j=1}^{N} U_j\, \varphi_j(x),
\]
where $\{\varphi_j\}$ is a basis of \concept{tent functions}
(piecewise polynomial, supported on a few mesh elements each). The
$U_j$ are the unknown \concept{degrees of freedom (DoFs)}. Plugging
this into the weak form turns each PDE into an equation per basis
function, giving a square nonlinear system $F(U) = 0$ for
$U = (U_1, \dots, U_N)$.

\paragraph{Why FEM rather than finite differences?} FEM handles
non-rectangular domains, variable mesh resolution, and Robin
boundary conditions natively. We do not need any of those
features in the strict sense (we use a 2D rectangular mesh with
graded refinement near the electrode), but FEM gives us:
(i) trivial mesh refinement near the EDL singular layer;
(ii) a well-developed adjoint capability via Firedrake's
\texttt{firedrake.adjoint} (Pyadjoint), which we will need for the
inverse problem when it un-pauses; and
(iii) symbolic assembly of the Jacobian for Newton, which would be
brittle to write by hand.

\paragraph{Firedrake.} The Python library that makes this all
declarative. You write the weak form with operators like
\texttt{inner(grad(u), grad(v))*dx}, declare a function space, and
ask it to \texttt{solve(F == 0, u)}. Under the hood it generates C
kernels for assembly (via TSFC), assembles the sparse Jacobian, and
hands it to PETSc for the linear solve. You will not need to read
the Firedrake internals; treat it as a tool.

\subsection{The mesh}

We use a 2D \concept{graded rectangle} of width $\sim 1$ scaled
unit and height $\hat{L} = 1$ (which corresponds to physical
$L_{\mathrm{eff}} \sim 16$--$100\,\mu$m). The graded grid puts most
nodes near the electrode at $x = 0$ where the EDL lives:
\[
  y_j = \hat{L} \cdot \frac{e^{\beta j/N_y} - 1}{e^{\beta} - 1},
  \qquad \beta \approx 3,
\]
giving $\sim 80$ nodes spaced over five decades of physical
distance. The mesh constructor lives in
\texttt{Forward/bv\_solver/mesh.py} (\texttt{make\_graded\_rectangle\_mesh}).
\warn{Keep $\hat{L}$ matched to whatever \texttt{l\_eff\_m} the
solver was built with}: an $\hat{L}$-mismatch silently puts the bulk
Dirichlet BC at the wrong physical location.

% =========================================================================
\section{Newton's method and where it goes wrong}

Newton's method (you have seen it):
\[
  J(U^k)\,\Delta U = -F(U^k), \qquad U^{k+1} = U^k + s_k\,\Delta U,
\]
with line-search step $s_k \in (0,1]$. Firedrake delegates the
linear solve to PETSc; we use a sparse direct solver (MUMPS) at
this problem size since the Jacobian fits in memory.

\subsection{Why cold-start often fails}

Two failure modes dominate:

\begin{enumerate}[leftmargin=*]
\item \textbf{The Jacobian is poorly conditioned} when $|\eta|$ is
      large. The BV exponentials (doc 01) saturate
      to $10^{40}$-ish numbers; pre-clipping helps but does not
      eliminate the problem. We use $\texttt{exponent\_clip} = 100$
      (clip on $\eta_{\mathrm{scaled}} = (V - \Eeq)/V_T$ before
      multiplication by $\alpha n_e$); see
      \texttt{docs/solver/clipping\_conventions.md}.
\item \textbf{The initial guess is too far from the solution.} For
      a target like $\Vrhe = +1.0$ V at production $\kzero$, a
      cold-start from a uniform-bulk profile lands the Newton
      iteration on the wrong side of a manifold and diverges.
\end{enumerate}

\subsection{Continuation: walk to the solution one step at a time}

The fix is \concept{continuation}: replace the hard problem with a
sequence of easier problems whose solutions are close together.
We use three layered continuations:

\begin{description}[leftmargin=*]
  \item[Voltage continuation.] Solve at a sequence of $\Vrhe$
      values, using each converged solution as the warm-start
      initial guess for the next $\Vrhe$. Implemented as the
      ``warm-walk'' phase of the orchestrator.
  \item[$\kzero$ continuation.] Within a single $\Vrhe$, scale the
      kinetic prefactor down by many orders of magnitude
      ($10^{-12}, 10^{-9}, 10^{-6}, 10^{-3}, 10^{0}$ relative to
      target), solve, then scale back up. Tagged
      \texttt{initial\_scales} in the code.
  \item[Phase 6$\beta$-only: $\lambda_{\mathrm{hyd}}$ and
      $k_w^{\mathrm{eff}}$ ladders.] Smooth turn-on of the cation
      hydrolysis source ($\lambda$) and water self-ionization rate
      ($k_w$) via outer loops, similar to a homotopy from the
      Phase 5$\gamma$ solution to the Phase 6$\beta$ target.
\end{description}

The orchestrator that owns these loops is in
\texttt{Forward/bv\_solver/anchor\_continuation.py}, specifically:
\begin{itemize}[leftmargin=*]
\item \texttt{solve\_anchor\_with\_continuation}: build a converged
      ``anchor'' solution at one well-chosen voltage, ramping
      $\kzero$ and $k_w^{\mathrm{eff}}$ ladders to target.
\item \texttt{extract\_preconverged\_anchor}: repackage the anchor
      for warm-start use.
\item \texttt{solve\_grid\_with\_anchor}: warm-walk the converged
      anchor across a list of voltages.
\end{itemize}
This pair is the production entry point for the multi-ion + Stern
+ Phase 6$\alpha$/6$\beta$ stack. The older
\texttt{solve\_grid\_per\_voltage\_cold\_with\_warm\_fallback}
(``C+D'' orchestrator) is used for the legacy single-counter-ion
ClO$_4^-$ stack but cold-fails on the multi-ion stack near
$\Vrhe \approx +0.55$ V (13/13 fail). \warn{Do not pick C+D for
deck-aligned work}; use the anchor-and-grid pair.

\subsection{Picard for the outer loop}

Newton converges quadratically near the solution but needs a smooth
residual. The cation hydrolysis coverage equation
(doc 02, $\dot\Gamma$ formula) has algebraic clamps and is not smooth
enough for Newton to like it. We solve it with a separate outer
\concept{Picard iteration}: solve the inner PNP+BV Newton with
$\Gamma$ held fixed; update $\Gamma$ from the new fields; repeat
until $\Gamma$ stops moving. The Picard helper lives in
\texttt{Forward/bv\_solver/cation\_hydrolysis.py}.

\subsection{Adaptive ladders (step 9.5 onwards)}

Late Phase 6$\beta$ added \texttt{AdaptiveLadder} in the
\texttt{Forward/bv\_solver/anchor\_continuation.py}, which
arithmetic-bisects the $\lambda_{\mathrm{hyd}}$ ramp when a Newton
solve fails partway up. The
\texttt{warm\_start\_floor} option in step 9.5 (May 2026) added
the ability to bisect from the warm-start floor ($\lambda = 0$),
rescuing 3/14 high-$k_{\mathrm{hyd}}$ configurations in the step 9
B.2 sweep.

% =========================================================================
\section{Repository layout}

A 30-second tour of the top of \texttt{PNPInverse/}:

\begin{center}
\begin{tabular}{p{0.30\linewidth} p{0.66\linewidth}}
\toprule
\textbf{Path} & \textbf{What lives there} \\
\midrule
\texttt{Forward/} & Forward solver code. The active surface. \\
\texttt{Forward/bv\_solver/} & The PNP--BV Firedrake package. Most of your work will be here. \\
\texttt{scripts/studies/} & Driver scripts. One file per study. \\
\texttt{scripts/derive/} & Data-derivation scripts (Tafel slopes, etc.) \\
\texttt{scripts/verification/} & MMS and other verification tests. \\
\texttt{calibration/} & Top-level constants (\texttt{v10b.py}). Firedrake-free. \\
\texttt{Nondim/} & Physical constants \& nondim transforms. \\
\texttt{tests/} & Pytest tests. \texttt{slow}-marked ones need Firedrake. \\
\texttt{docs/} & All written documentation. \\
\texttt{StudyResults/} & Output of every run we have kept. JSONs + plots. \\
\texttt{data/} & \emph{Gitignored.} Experimental data (273 MB). Ask for the drop. \\
\texttt{Inverse/}, \texttt{FluxCurve/}, \texttt{Surrogate/} & \emph{Paused.} Don't touch. \\
\texttt{archive/} & Old code, reference only. \\
\texttt{CLAUDE.md} & The hard rules. Read once, refer often. \\
\texttt{README.md} & The project narrative. \\
\bottomrule
\end{tabular}
\end{center}

\subsection{Inside \texttt{Forward/bv\_solver/}}

\begin{center}
\begin{tabular}{p{0.32\linewidth} p{0.62\linewidth}}
\toprule
\textbf{File} & \textbf{Role} \\
\midrule
\texttt{forms\_logc.py}        & Residual builder, log-c on all species (legacy backend, kept for cross-check). \\
\texttt{forms\_logc\_muh.py}   & Residual builder, $\mu_H$ on protons (\textbf{production backend}). \\
\texttt{boltzmann.py}          & Bikerman analytic counter-ion expressions (multi-ion). \\
\texttt{water\_ionization.py}  & Phase 6$\alpha$ $\HtwoO \rightleftharpoons \Hp + \OHm$ closure. \\
\texttt{cation\_hydrolysis.py} & Phase 6$\beta$ $\Gamma_{\mathrm{MOH}}$ Picard + Singh pK$_a$ shift. \\
\texttt{anchor\_continuation.py} & Phase 5$\gamma$ orchestrator: anchor + grid + ladders. \\
\texttt{grid\_per\_voltage.py} & Legacy C+D orchestrator (ClO$_4^-$ stack only). \\
\texttt{picard\_ic.py}         & Shared Picard outer loop for the \texttt{debye\_boltzmann} IC. \\
\texttt{dispatch.py}           & Routes \texttt{build\_context}, \texttt{build\_forms},
                                  \texttt{set\_initial\_conditions} to the right backend. \\
\texttt{config.py}             & Reads/validates the \texttt{bv\_convergence} options. \\
\texttt{diagnostics.py}        & Per-voltage diagnostics (BV-rate components, $F_0$, $\theta$, \dots). \\
\texttt{rrde\_observables.py}  & Postprocesses to ring/disk currents, selectivity, $n_e$. \\
\texttt{validation.py}         & W1--W5 sanity checks on a converged solution. \\
\texttt{mesh.py}               & Graded rectangle mesh constructor. \\
\bottomrule
\end{tabular}
\end{center}

\subsection{Inside \texttt{scripts/studies/}}

Each driver corresponds to one named study or one paper-figure
target. The current ones to know:

\begin{center}
\begin{tabular}{p{0.42\linewidth} p{0.52\linewidth}}
\toprule
\textbf{Driver} & \textbf{What it does} \\
\midrule
\texttt{phase6b\_step6\_plumbing\_ablation.py}
  & K$^+$/SO$_4^{2-}$ + cation hydrolysis at $V_{\mathrm{kin}}$ with
    Step 6 ablation flags. \\
\texttt{phase6b\_v10a\_phase\_A2\_v\_kin.py}
  & Phase A.2 baseline reproduction at $V_{\mathrm{kin}}$ with v10a
    Langmuir cap. \\
\texttt{phase6b\_v10a\_v\_sweep\_diagnostic.py}
  & V-sweep diagnostic for $V_{\mathrm{kin}}$ selection. \\
\texttt{l\_eff\_transport\_sweep\_csplus\_so4.py}
  & Cs$^+$/SO$_4^{2-}$ multi-ion + Phase 6$\alpha$ validation. \\
\texttt{mangan\_full\_grid\_csplus\_so4.py}
  & Cs$^+$/SO$_4^{2-}$ deck page-15 V-band sweep. \\
\bottomrule
\end{tabular}
\end{center}

\subsection{Where the tests live}

\texttt{tests/} contains pytest regression tests. Two markers:
\texttt{not slow} (no Firedrake required) and \texttt{slow} (need
the Firedrake venv active). Run them with
\begin{lstlisting}
python -m pytest -m "not slow"
python -m pytest -m slow
\end{lstlisting}
Phase 6$\beta$ has its own test files
(\texttt{test\_phase6b\_v9\_*.py},
\texttt{test\_phase6b\_step10\_phase\_D\_*.py}).

% =========================================================================
\section{Anatomy of a study run}

\subsection{Outputs}

A study run drops a folder under \texttt{StudyResults/<study\_name>/}
containing:

\begin{center}
\begin{tabular}{p{0.30\linewidth} p{0.64\linewidth}}
\toprule
\textbf{File} & \textbf{Contents} \\
\midrule
\texttt{iv\_curve.json}     & The current--voltage curve at each grid voltage; the headline observable. \\
\texttt{summary.json}       & Per-voltage convergence flags, BV-rate components, surface concentrations, $\Gamma$. \\
\texttt{verdict.json}       & PASS/FAIL of the study's predefined acceptance gates. \\
\texttt{*.png}              & Plots (cd, peroxide, surface pH, etc.). \\
\texttt{run.log}            & Console output of the driver. \\
\bottomrule
\end{tabular}
\end{center}

When you wonder ``did anyone already run this?'', open the
\texttt{StudyResults/} subdirectory and read the existing
\texttt{summary.json}. \warn{A full sweep can take minutes to
hours} (production = $\sim 30$ min for 8 combos $\times$ 13
voltages); do not regenerate without checking.

\subsection{The convergence flags}

Each voltage point in \texttt{summary.json} has a
\texttt{converged} flag and an \texttt{exit\_status}. The most
common values:
\begin{itemize}[leftmargin=*]
\item \texttt{newton\_converged}: rosy, residual $L^2$ below tol.
\item \texttt{ladder\_partial}: a continuation ramp made it part of
      the way; the recorded fields are at a partially-converged
      $\kzero$ scale.
\item \texttt{cold\_fail}: cold-start Newton diverged; warm-fallback
      may or may not have rescued.
\item \texttt{stagnant}: Newton stopped reducing residual.
\end{itemize}

\subsection{Cross-checking: validation pass}

\texttt{validation.py} runs a fast set of W1--W5 sanity checks on
each converged field set. They check things like clip saturation
(W1), counter-ion depletion (W5), BV-rate vs.\ Faraday's law
consistency. These are cheap; always look at them before believing
a result.

% =========================================================================
\section{Writing a new diagnostic}

The most common first contribution from a new collaborator is
adding a diagnostic plot or a new observable. The pattern:

\begin{enumerate}[leftmargin=*]
\item Pick the quantity you want to compute (e.g.\ ``volume-averaged
      pH in the diffusion layer'').
\item Write a function in
      \texttt{Forward/bv\_solver/diagnostics.py} that takes
      \texttt{ctx} (the solver context, a dict of UFL functions)
      and returns a Python float.
\item Add a regression test under \texttt{tests/} that pins it on a
      reference field.
\item Wire the diagnostic into the driver script you care about.
\item Run the study; verify the new field shows up in
      \texttt{summary.json}.
\end{enumerate}

The repo has dozens of examples of this pattern;
\texttt{diagnostics.py} itself is a good model.

% =========================================================================
\section{The adjoint pipeline (paused, but useful background)}

\texttt{firedrake.adjoint} (Pyadjoint) provides reverse-mode
automatic differentiation through Firedrake solves. When the inverse
work resumes, this is what gives us $\partial L/\partial \theta$
where $L$ is a misfit and $\theta = \{\kzero, \alpha, \dots\}$. The
hard rule for adjoints is:
\begin{callout}{Adjoint tape hygiene}
Any unannotated cold-ramp or continuation work must be wrapped in
\texttt{with adj.stop\_annotating(): ...}. Otherwise the tape
records intermediate ladder solves and the gradient is wrong /
slow. This is in CLAUDE.md and in the inverse handoff
\texttt{docs/inverse/CHATGPT\_HANDOFF\_10\_LM\_TIKHONOV\_BASIN\_GEOMETRY.md}.
\end{callout}

You will not need to write adjoints in your first months. Just know
that ``annotate'' and ``stop\_annotating'' are jargon for ``record
this solve onto the tape so we can backprop through it later''.

% =========================================================================
\section{Hard rules summary}

These are the project's CLAUDE.md hard rules, condensed:

\begin{enumerate}[leftmargin=*]
\item \textbf{Use anchor-and-grid (or C+D for legacy ClO$_4^-$),
      not Strategy B.} Strategy B
      (\texttt{solve\_grid\_with\_charge\_continuation}) Phase-1
      cold-fails on the logc + counter-ion stack.
\item \textbf{\texttt{exponent\_clip = 100} is the only
      PC-trustworthy setting.} Lower clips create fictitious
      peroxide currents at high $\Vrhe$.
\item \textbf{The IC and the residual must agree about steric
      saturation.} Bikerman IC + ideal residual cold-fails; same
      goes vice versa.
\item \textbf{Use the parallel R$_{2e}$ + R$_{4e}$ topology}
      (Ruggiero 2022). Sequential R$_0$/R$_1$ is wrong and is kept
      only for backward-compat.
\item \textbf{Check \texttt{data/EChem Reactor Modeling-Seitz-Mangan/}
      first.} Do not conjecture about the experimental physics.
      The hypothesis is cation hydrolysis at the polarized OHP, not
      sulfate, not carbonate, not water self-ionization.
\item \textbf{$C_S = 0.20$ F/m$^2$ is the literature-anchored
      target} (Bohra--Koper--Choi consensus). Lower values were
      used historically without citation.
\end{enumerate}

\onboardingfooter
\end{document}
